
\documentclass[10pt,letterpaper]{article}

% ===== Tipografía y codificación (pdfLaTeX) =====
\usepackage[utf8]{inputenc}      % Entrada en UTF-8
\usepackage[T1]{fontenc}         % Salida con codificación T1
\usepackage{lmodern}             % Fuente moderna

                       % separador decimal con punto si lo deseas

% ===== Márgenes =====
\usepackage[margin=2cm]{geometry}     % ajusta a tu gusto

% ===== Matemáticas =====
\usepackage{amsmath, amssymb, amsthm}

% ===== Gráficos e imágenes =====
\usepackage{graphicx}                 % reemplaza epsfig
\usepackage{float}                    % [H] colocación estricta si la necesitas
\usepackage{booktabs}                 % tablas bonitas
\usepackage{xcolor}                   % colores


% ===== Columnas múltiples (bloques en dos columnas) =====
\usepackage{multicol}
\setlength{\columnsep}{14pt}

\renewcommand{\theenumi}{\arabic{enumi}$^{\circ}$}
\renewcommand{\labelenumi}{\theenumi.}




% ===== Configuración general =====
\setlength{\parindent}{0pt}
\setlength{\parskip}{4pt}
\linespread{1.05}
%%%%%%%%%%%%%%%%%%%%%%%%%%%%%%%%%%%%%%%%%%%%%%%%%%%%%%%%%%%
\begin{document}
	\pagestyle{empty}
	

\textbf{PONTIFICIA UNIVERSIDAD JAVERIANA CALI}

Departamento de Ciencias Naturales y Matemáticas

300MAE013 - A - Teoría de Probabilidades

\begin{flushright}
Profesor: Daniel Enrique González Gómez \\  
Cali, septiembre 10 de 2026
\end{flushright}

\begin{center}
{\bf EXAMEN FINAL}
\end{center}

		\begin{tabular}{|p{1.5cm}|p{8.2cm}|p{2.2cm}|p{2.2cm}|p{2.4cm}|}
	\hline
	&&&&\\
	Habilidad & Indicador   &P1. & P2. &  NOTA\\  
	&&&&\\ \hline 
	(50\%) A& Habilidad para aplicar conocimientos de matemáticas, ciencias e ingeniería    &&&   \\ 
	&&&& \\ \hline              
	(50\%) G& Habilidad para comunicarse efectivamente &&&\\  
	&&&& \\  
	\hline  
\end{tabular}

\begin{tabular}{|p{18.2cm}|}
\hline      
{\textbf{OBSERVACIONES}:} \\ \\ 
\\ \\ 
\\ \\ 
\\ \\ 

\hline

\end{tabular}

\begin{multicols}{2}
%%%%%%%%%%%%%%%%%%%%%%%%%%%%%%%%%%%%%%%%%%%%%%%%%%%%%%%%%%%%%%%%%%	
\begin{itemize}
\item[1.] 



%%%%%%%%%%%%%%%%%%%%%%%%%%%%%%%%%%%%%%%%%%%%%%%%%%%%%%%%%%%%%%%%%
\item[2.] En una empresa de distribución, se están evaluando tres tipos de estrategias de capacitación implementadas para el personal comercial: capacitación presencial tradicional, capacitación virtual interactiva y capacitación basada en acompañamiento en campo. Se ha determinado que el 25\% de los vendedores recibió capacitación presencial tradicional, el 50\% capacitación virtual interactiva, y el 25\% capacitación basada en acompañamiento en campo.

Además, se ha observado que la probabilidad de que un vendedor mantenga un desempeño satisfactorio después de 5 años es del 80\% para quienes recibieron capacitación presencial tradicional, del 60\% para quienes recibieron capacitación virtual interactiva, y del 40\% para quienes recibieron capacitación basada en acompañamiento en campo.

Si se selecciona al azar un vendedor que presenta desempeño satisfactorio después de 5 años, ¿cuál es la probabilidad de que haya recibido capacitación presencial tradicional?

%%%%%%%%%%%%%%%%%%%%%%%%%%%%%%%%%%%%%%%%%%%%%%%%%%%%%%%%%%%%%%%%%
\end{itemize}
%%%%%%%%%%%%%%%%%%%%%%%%%%%%%%%%%%%%%%%%%%%%%%%%%%%%%%%%%%%%%%%%%
\end{multicols}
%%%%%%%%%%%%%%%%%%%%%%%%%%%%%%%%%%%%%%%%%%%%%%%%%%%%%%%%%%%%%%%%%
\end{document}
